\section*{Capítulo \ref{ch:g6:density-floating} --- La densidad: por qué flotan las cosas}

\begin{solution}{exo:g6:density-floating:1}
La \omterm{def:g3:measuring-mass:mass}{masa} de un \omterm{def:g6:mass-and-volume:volume}{centímetro cúbico} de la sustancia, en
\omterm{def:g3:measuring-mass:units}{gramos}; la del agua es \qty{1}{g} por \unit{cm^3}. La regla: lo menos
denso que el agua \omterm{def:g1:floating-and-sinking:words}{flota} y lo más denso \omterm{def:g1:floating-and-sinking:words}{se hunde}.
\end{solution}

\begin{solution}{exo:g6:density-floating:2}
El roble \omterm{def:g1:floating-and-sinking:words}{flota}; la piedra \omterm{def:g1:floating-and-sinking:words}{se hunde}; el corcho \omterm{def:g1:floating-and-sinking:words}{flota}; el
hierro \omterm{def:g1:floating-and-sinking:words}{se hunde}; el hielo \omterm{def:g1:floating-and-sinking:words}{flota} --- por los pelos.
\end{solution}

\begin{solution}{exo:g6:density-floating:3}
$120 \div 80 = 1.5$: \omterm{def:g6:density-floating:density}{densidad} $1.5$ --- más denso que el agua, así
que \omterm{def:g1:floating-and-sinking:words}{se hunde}.
\end{solution}

\begin{solution}{exo:g6:density-floating:4}
$158 \div 20 = 7.9$: la firma del hierro.
\end{solution}

\begin{solution}{exo:g6:density-floating:5}
La \omterm{def:g6:density-floating:density}{densidad} pertenece a la sustancia, no al trozo: todos los
\omterm{def:g6:mass-and-volume:volume}{centímetros cúbicos} de roble llevan la misma ración de \omterm{def:g3:measuring-mass:units}{gramos},
sean de una astilla o de una viga. Eso convierte a la \omterm{def:g6:density-floating:density}{densidad} en
un carné de identidad para sustancias desconocidas.
\end{solution}

\begin{solution}{exo:g6:density-floating:6}
Los \omterm{def:g3:solids-liquids-gases:liquid}{líquidos} se apilan por \omterm{def:g6:density-floating:density}{densidad}, con el más denso abajo:
la miel $1.4$ debajo del agua $1.0$ y esta debajo del aceite $0.9$. La
uva ($1.05$) \omterm{def:g1:floating-and-sinking:words}{se hunde} por el aceite y por el agua y se posa en el
suelo de miel --- más densa que los dos pisos de arriba y menos densa
que la miel.
\end{solution}

\begin{solution}{exo:g6:density-floating:7}
Casi todos los \omterm{def:g3:solids-liquids-gases:solid}{sólidos} son más densos que su propio \omterm{def:g3:solids-liquids-gases:liquid}{líquido} y
se hunden en él. El hielo \omterm{def:g1:floating-and-sinking:words}{flota} porque el agua se hincha al
\omterm{def:g2:ice-water-steam:meltfreeze}{congelarse} --- la misma sustancia en una décima parte más de sitio da
una \omterm{def:g6:density-floating:density}{densidad} de $0.9 < 1$.
\end{solution}

\begin{solution}{exo:g6:density-floating:8}
La madera firma unos $0.6$: todos los \omterm{def:g6:mass-and-volume:volume}{centímetros cúbicos} del
tronco pujan por debajo del \qty{1}{g} del agua, así que el tronco
\omterm{def:g1:floating-and-sinking:words}{flota} por enorme que sea. El bronce firma cerca de $9$: todos los
\omterm{def:g6:mass-and-volume:volume}{centímetros cúbicos} de la moneda pujan por encima del agua, así que
\omterm{def:g1:floating-and-sinking:words}{se hunde} por pequeña que sea. El tamaño no importó nunca; la
\omterm{def:g3:measuring-mass:mass}{masa} por sitio ocupado importó siempre.
\end{solution}

\begin{solution}{exo:g6:density-floating:9}
$\qty{1}{L} = \qty{1000}{cm^3}$, así que $800 \div 1000 = 0.8$
\omterm{def:g3:measuring-mass:units}{gramos} por \omterm{def:g6:mass-and-volume:volume}{centímetro cúbico}. El hielo, con $0.9$, es
\emph{más denso} que este \omterm{def:g3:solids-liquids-gases:liquid}{líquido}: el cubito \omterm{def:g1:floating-and-sinking:words}{se hunde} en él.
\end{solution}

\begin{solution}{exo:g6:density-floating:10}
El paquete que \omterm{def:g1:floating-and-sinking:words}{flota} es el casco entero: la cáscara de acero
\emph{más} el enorme \omterm{def:g6:mass-and-volume:volume}{volumen} de \omterm{def:g2:air-around-us:air}{aire} que encierra. Promediada
sobre todo ese sitio, la \omterm{def:g6:density-floating:density}{densidad} del paquete cae por debajo de $1$,
y la regla dice \omterm{def:g1:floating-and-sinking:words}{flotar}. Una vía de agua deja que el agua sustituya
al \omterm{def:g2:air-around-us:air}{aire}: la \omterm{def:g6:density-floating:density}{densidad} del paquete sube hacia la del propio
acero, pasa de $1$ --- y la regla, impasible, dice \omterm{def:g1:floating-and-sinking:words}{hundirse}.
\end{solution}

\begin{solution}{exo:g6:density-floating:11}
El árbitro ha cambiado: frente al $1.03$ del agua de mar, el cuerpo de
un nadador, cercano a $1$, queda cómodamente del lado de la flotación,
así que el mar te lleva más alto. Al entrar en agua dulce ($1.0$), el
barco pierde ese margen y \omterm{def:g1:floating-and-sinking:words}{se hunde} más: su línea de flotación sube por
el casco.
\end{solution}

\begin{solution}{exo:g6:density-floating:12}
Medio sumergida, la esfera \omterm{def:g1:floating-and-sinking:words}{se hunde} en el agua exactamente lo que
debería \omterm{def:g1:floating-and-sinking:words}{hundirse} un paquete con la mitad de la \omterm{def:g6:density-floating:density}{densidad} del agua:
la cáscara más el \omterm{def:g2:air-around-us:air}{aire} promedian alrededor de $0.5$ \omterm{def:g3:measuring-mass:units}{gramos} por
\omterm{def:g6:mass-and-volume:volume}{centímetro cúbico} --- y quien más aligera es el \omterm{def:g2:air-around-us:air}{aire} escondido.
Según entra agua, esta sustituye al \omterm{def:g2:air-around-us:air}{aire}, la \omterm{def:g6:density-floating:density}{densidad} media del
paquete sube y la esfera \omterm{def:g1:floating-and-sinking:words}{se hunde} cada vez más; en el momento en que
esa media pasa del $1.0$ del agua, la regla de la flotación
cambia de bando y la esfera se va al fondo.
\end{solution}

\begin{solution}{pb:g6:density-floating:1}
\textbf{1.} Nada: cambiar oro por una \emph{\omterm{def:g3:measuring-mass:mass}{masa}} igual de plata
mantiene contenta a la \omterm{def:g3:levers-and-scales:scale}{balanza} --- la \omterm{def:g3:measuring-mass:mass}{masa} por sí sola no puede ver
la sustitución.
\textbf{2.} $1930 \div 19.3 = \qty{100}{cm^3}$.
\textbf{3.} \qty{130}{cm^3} --- lo desbordado es el volumen de la
corona, por desplazamiento.
\textbf{4.} La corona ocupa \qty{30}{cm^3} más de sitio del que
debería ocupar el oro puro de su \omterm{def:g3:measuring-mass:mass}{masa}: dentro se esconde algo más
voluminoso que el oro.
\textbf{5.} $1930 \div 130 \approx 14.8$.
\textbf{6.} Oro $19.3$ --- corona $14.8$ --- plata $10.5$: la corona
cae entre las dos firmas.
\textbf{7.} ``Majestad, cada \omterm{def:g6:mass-and-volume:volume}{centímetro cúbico} de oro puro lleva
\qty{19.3}{g}; vuestra corona lleva solo unos $15$ --- de oro puro no
es''.
\textbf{8.} El intervalo da \omterm{def:g6:density-floating:density}{densidades} desde
$1900 \div 130 \approx 14.6$ hasta $1960 \div 130 \approx 15.1$ ---
ni de lejos $19.3$: ningún error verosímil de la \omterm{def:g3:levers-and-scales:scale}{balanza} salva al
orfebre.
\textbf{9.} La plata es menos densa: cada \omterm{def:g3:measuring-mass:units}{gramo} de oro robado
devolvió un \omterm{def:g3:measuring-mass:units}{gramo} de plata que necesita \emph{más sitio} --- la
sustitución \omterm{def:g3:measuring-mass:units}{gramo} a \omterm{def:g3:measuring-mass:units}{gramo} conserva la \omterm{def:g3:measuring-mass:mass}{masa} e infla el volumen.
\textbf{10.} Parte de oro: $965 \div 19.3 = \qty{50.0}{cm^3}$; parte
de plata: $965 \div 10.5 \approx \qty{91.9}{cm^3}$.
\textbf{11.} Total $\approx \qty{142}{cm^3}$ --- más que los $130$
medidos: la conjetura mitad y mitad usó demasiada plata; el robo
verdadero fue menor (el volumen medido queda entre $100$ y $142$).
\textbf{12.} ``Probado: la corona no es de oro puro. Método: su
\omterm{def:g3:measuring-mass:mass}{masa} repartida entre su \omterm{def:g6:mass-and-volume:volume}{volumen} medido con agua da una
\omterm{def:g6:density-floating:density}{densidad} muy por debajo de la firma inmutable del oro. Y la corona
no necesitó ni un arañazo: el agua leyó su secreto desde fuera --- ese
es el triunfo''.
\end{solution}
